\documentclass[onecolumn]{revtex4}
\usepackage{ragged2e}
\usepackage{amsmath}
\usepackage{graphicx}
\usepackage{bm}
\usepackage{soul} % for highlight  \hl{...}

\usepackage[usenames,dvipsnames]{color}

\definecolor{darkred}{RGB}{0,0,196}
%\usepackage[colorlinks=true,linkcolor=darkred,citecolor=darkred,urlcolor=darkred]{hyperref}
% Adding or removing % in the last \usepackage can change the styles of refs

\renewcommand\baselinestretch{1.265} %1.265
% When the line spacing and figure size are not appropriate
% and result in compilation failure, they can be modified

\begin{document}

\title{Extraction of Effective Temperatures from Tsallis Fits to
Transverse-Momentum Spectra of Direct Photons and Heavy Quarkonia in
High-Energy Collisions} \vspace{0.5cm}

\author{Peng-Cheng~Zhang$^{1,}$\footnote{202312602003@email.sxu.edu.cn},
Hailong~Zhu$^{1,}$\footnote{Correspondence: zhuhl@sxu.edu.cn} and
Fu-Hu~Liu$^{1,}$\footnote{Correspondence: fuhuliu@163.com;
fuhuliu@sxu.edu.cn}}

\affiliation{$^1$Institute of Theoretical Physics \& College of
Physics and Electronic Engineering, State Key Laboratory of
Quantum Optics Technologies and Devices \& Collaborative
Innovation Center of Extreme Optics, Shanxi University, Taiyuan
030006, China}

\begin{abstract}

\vspace{0.5cm}

\noindent {\bf Abstract:} 
\\
\\
{\bf Keywords:} Tsallis statistics; effective temperature
\\
\\
{\bf PACS numbers:}
\\
\end{abstract}

%\date{\today}
\maketitle

\parindent=15pt


\section{Results and Discussion}

A central feature of the nonprompt direct-photon spectra~\cite{PHENIX2024} 
is that they are not described by a single
exponential: the local inverse slope rises with transverse momentum, and
the data were accordingly characterized by two separate exponential fits,
yielding $T_{\rm eff} \approx 0.26$~GeV/$c$ in the range $0.8 < p_T < 1.9$~GeV/$c$
and $T_{\rm eff} \approx 0.38$~GeV/$c$ in $2 < p_T < 4$~GeV/$c$. The present
analysis shows that this evolving slope is captured naturally, and without
piecewise treatment, by a single Tsallis form fitted over
the full momentum range.

The equivalence follows directly from the structure of the Tsallis
function. For $dN/dp_T \propto p_T\,[1+(q-1)\,p_T/T]^{q/(1-q)}$, the local
inverse slope is not constant but a smooth, monotonically increasing
function of $p_T$,
\begin{equation}
    T_{\rm eff}(p_T) \equiv
    -\left(\frac{d\ln(dN/dp_T)}{dp_T}\right)^{-1},
\end{equation}
which reduces to the constant $T$ only in the Boltzmann limit $q\to 1$. A
value $q>1$ therefore encodes precisely the spectral hardening that the
two-exponential description represents through a discrete jump in
$T_{\rm eff}$. The single parameter $q$ thus subsumes the two independent
slopes of the piecewise fit. %这段添加到第二节公式里

This is borne out quantitatively. Evaluating $T_{\rm eff}(p_T)$ for the
0\%--20\% best-fit parameters ($q = 1.084$, $T = 0.128$~GeV) gives
$T_{\rm eff} \approx 0.26$~GeV/$c$ near $p_T = 1.3$~GeV/$c$ and
$T_{\rm eff} \approx 0.40$~GeV/$c$ near $p_T = 3$~GeV/$c$, in close
agreement with the two inverse slopes extracted independently~\cite{PHENIX2024} for the corresponding $p_T$ intervals. The two
regimes are recovered not as separate fits but as the low- and high-$p_T$
limits of one continuous curve.

The Tsallis description is therefore both more economical and more
physically suggestive. It replaces two phenomenological slopes by a single
functional form with one shape parameter $q$, and it admits a transparent
interpretation within superstatistics: a Tsallis spectrum arises from a
superposition of Boltzmann emitters with a distribution of inverse
temperatures $\beta$, with $q-1 \simeq \mathrm{Var}(\beta)/\langle\beta\rangle^2$.
The hardening described piecewise by two slopes is thus reframed as a
single spectrum emitted by a system spanning a range of temperatures---a
picture consistent with photon emission integrated over the cooling history
of the expanding fireball.


\begin{figure*}[htb!]
\begin{center}
\includegraphics[width=15.5cm]{/home/weyl/codexdem/Au+Au/nonprompt/cenm1_nonprompt.pdf}
\end{center}
\justifying\noindent {Figure 1a. Nonprompt direct-photon spectra in Au$+$Au collisions at
    $\sqrt{s_{NN}}=200$~GeV for the four centrality classes, with the Tsallis fits
    (curves) evaluated at the best-fit $(q,T,\mathrm{Nor})$ values.}
\end{figure*}




\begin{figure*}[htb!]
\begin{center}
\includegraphics[width=15.5cm]{/home/weyl/codexdem/Au+Au/nonprompt/qt_contours_err_figure16_0_10.pdf}
\end{center}
\justifying\noindent {Figure 1d. Tsallis-fit confidence regions in the $(q, T)$ plane for 
    four centrality classes. Closed curves show the 68.3\% joint confidence 
    region ($\Delta\chi^2=2.30$); filled circes mark the best fits, with bars 
    giving single-parameter $1\sigma$ errors ($\Delta\chi^2=1$).}
\end{figure*}



{\small
\begin{thebibliography}{99}

\setlength{\itemsep}{2pt}

\bibitem{PHENIX2024}%[1]
N. J. Abdulameer {\it et al.} (PHENIX Collaboration), ``Nonprompt
direct-photon production in Au$+$Au collisions at $\sqrt{s_{NN}}=200$~GeV,"
{\it Physical Review C} 109, no. 4 (2024): 044912,
https://doi.org/10.1103/PhysRevC.109.044912.


\end{thebibliography}}


\end{document}
